\documentclass[10pt]{extarticle}
\usepackage[a4paper,margin=1.5cm]{geometry}
\usepackage[T1]{fontenc}
\usepackage{amsmath,amssymb}
\usepackage{graphicx}
\usepackage[export]{adjustbox}
\usepackage{multicol}
\usepackage{enumitem}
\usepackage{booktabs,array,tabularx}
\setlength{\columnsep}{0.7cm}
\setlength{\parindent}{0pt}
\setlength{\parskip}{3pt}
\pagestyle{empty}
\makeatletter
% figure registry (figdefs.tex is generated by the extract step)
\newcommand{\qfigdef}[4]{\@namedef{qf@file@#1}{#2}\@namedef{qf@w@#1}{#3}\@namedef{qf@wide@#1}{#4}}
\newcommand{\qf@check}[1]{\@ifundefined{qf@file@#1}{\PackageError{score2}{Unknown figure id #1}{}}{}}
\newcommand{\qf@img}[4][]{\includegraphics[width=#2\dimexpr\@nameuse{qf@w@#4}pt\relax,max width=#3,max height=0.45\textheight,#1]{\@nameuse{qf@file@#4}}}
% Full-width questions: finalize wraps every question that contains a wide figure in \qfullbegin...\qfullend,
% so the WHOLE question (stem, figures, options) is set across the page in reading order, then the two columns
% resume. The question's head - its number, stem and first figure (\qfullsplit marks the end of the head) -
% is kept on one page, so a number is never left at a page foot with its figure overleaf; the rest may break
% across pages normally, so no big gaps are left. Without a \qfullsplit, the whole question is kept together
% if it is under 45% of a page. Never write these yourself.
\newif\ifqf@full
\newif\ifqf@boxed
\newsavebox{\qf@box}
\newcommand{\qf@open}{\global\qf@fulltrue\par\global\qf@boxedtrue
  \setbox\qf@box\vbox\bgroup\hsize\textwidth\linewidth\textwidth}
\newcommand{\qf@emit}[1]{\par\egroup\global\qf@boxedfalse
  \ifdim\dimexpr\ht\qf@box+\dp\qf@box\relax<#1\textheight\noindent\usebox{\qf@box}\par\else\unvbox\qf@box\fi}
\newcommand{\qfullsplit}{\ifqf@boxed\qf@emit{0.85}\fi}
\newcommand{\qf@close}{\ifqf@boxed\qf@emit{0.45}\fi\global\qf@fullfalse}
\newcommand{\qfullbegin}{\end{multicols}\qf@open}
\newcommand{\qfullend}{\qf@close\begin{multicols}{2}}
\newcommand{\qfullstart}{\qf@open}
\newcommand{\qfulldone}{\qf@close}
% \qfig[scale]{id}: block figure, centred, directly where it is written (size set by the tool).
% Fallback if a wide figure is not inside a full-width question: break out of the columns in place - no floats.
\newcommand{\qfig}[2][1]{\qf@check{#2}%
  \ifqf@full\qf@blk{#1}{#2}\else
  \ifnum\@nameuse{qf@wide@#2}=1 \expandafter\@firstoftwo\else\expandafter\@secondoftwo\fi
  {\end{multicols}\begin{center}\qf@img{#1}{\textwidth}{#2}\end{center}\begin{multicols}{2}}%
  {\qf@blk{#1}{#2}}\fi}
\newcommand{\qf@blk}[2]{\par\smallskip{\centering\qf@img{#1}{\linewidth}{#2}\par}\smallskip}
% \qfigi[scale]{id}: inline figure (inside option lists / tables / minipages); never breaks columns.
\newcommand{\qfigi}[2][1]{\qf@check{#2}\qf@img[valign=c]{#1}{\linewidth}{#2}}
\makeatother
% structure macros - lint reads these to check completeness and figure placement
\newcommand{\papersec}[1]{\par\bigskip{\centering\Large\bfseries PAPER-#1\par}\medskip}
\newcommand{\partsec}[2]{\par\medskip\noindent{\large\bfseries PART-#1: #2}\par\smallskip}
\newcommand{\sect}[1]{\par\smallskip\noindent\textbf{#1}\par}
\newcommand{\qn}[1]{\par\smallskip\noindent\textbf{#1.}~}
\newcommand{\paraq}[2]{\par\smallskip\noindent\textbf{Paragraph for Questions #1 and #2}\par}
\newenvironment{qoptions}{\begin{itemize}[leftmargin=1.8em,itemsep=1pt,topsep=2pt,label={}]}{\end{itemize}}
\newcommand{\opt}[2]{\item[(#1)] #2}
% figure registry (generated by score2_tools.py finalize)
\qfigdef{p1_phy_s1_q01_1}{TNAS - Test 04 (18 Jun 2023)-img/p1_phy_s1_q01_1.png}{161.3}{0}
\qfigdef{p1_phy_s1_q02_1}{TNAS - Test 04 (18 Jun 2023)-img/p1_phy_s1_q02_1.png}{147.5}{0}
\qfigdef{p1_phy_s1_q03_1}{TNAS - Test 04 (18 Jun 2023)-img/p1_phy_s1_q03_1.png}{87.4}{0}
\qfigdef{p1_phy_s1_q04_1}{TNAS - Test 04 (18 Jun 2023)-img/p1_phy_s1_q04_1.png}{134.7}{0}
\qfigdef{p1_phy_s2_q07_1}{TNAS - Test 04 (18 Jun 2023)-img/p1_phy_s2_q07_1.png}{108.3}{0}
\qfigdef{p1_phy_s2_q07_2}{TNAS - Test 04 (18 Jun 2023)-img/p1_phy_s2_q07_2.png}{108.3}{0}
\qfigdef{p1_phy_s2_q08_1}{TNAS - Test 04 (18 Jun 2023)-img/p1_phy_s2_q08_1.png}{169.9}{0}
\qfigdef{p1_phy_s2_q09_1}{TNAS - Test 04 (18 Jun 2023)-img/p1_phy_s2_q09_1.png}{188.8}{0}
\qfigdef{p1_phy_s2_q10_1}{TNAS - Test 04 (18 Jun 2023)-img/p1_phy_s2_q10_1.png}{145.4}{0}
\qfigdef{p1_phy_s2_q11_1}{TNAS - Test 04 (18 Jun 2023)-img/p1_phy_s2_q11_1.png}{209.7}{0}
\qfigdef{p1_phy_s2_q12_1}{TNAS - Test 04 (18 Jun 2023)-img/p1_phy_s2_q12_1.png}{99.0}{0}
\qfigdef{p1_phy_s3_q01_1}{TNAS - Test 04 (18 Jun 2023)-img/p1_phy_s3_q01_1.png}{71.5}{0}
\qfigdef{p1_phy_s3_q03_1}{TNAS - Test 04 (18 Jun 2023)-img/p1_phy_s3_q03_1.png}{110.8}{0}
\qfigdef{p1_phy_s3_q04_1}{TNAS - Test 04 (18 Jun 2023)-img/p1_phy_s3_q04_1.png}{106.6}{0}
\qfigdef{p1_phy_s3_q05_1}{TNAS - Test 04 (18 Jun 2023)-img/p1_phy_s3_q05_1.png}{80.1}{0}
\qfigdef{p1_phy_s3_q06_1}{TNAS - Test 04 (18 Jun 2023)-img/p1_phy_s3_q06_1.png}{115.1}{0}
\qfigdef{p1_chem_s1_q01_1}{TNAS - Test 04 (18 Jun 2023)-img/p1_chem_s1_q01_1.png}{202.8}{0}
\qfigdef{p1_chem_s1_q03_1}{TNAS - Test 04 (18 Jun 2023)-img/p1_chem_s1_q03_1.png}{86.5}{0}
\qfigdef{p1_chem_s1_q03_2}{TNAS - Test 04 (18 Jun 2023)-img/p1_chem_s1_q03_2.png}{151.1}{0}
\qfigdef{p1_chem_s1_q04_1}{TNAS - Test 04 (18 Jun 2023)-img/p1_chem_s1_q04_1.png}{130.5}{0}
\qfigdef{p1_chem_s2_q07_1}{TNAS - Test 04 (18 Jun 2023)-img/p1_chem_s2_q07_1.png}{47.5}{0}
\qfigdef{p1_chem_s2_q11_1}{TNAS - Test 04 (18 Jun 2023)-img/p1_chem_s2_q11_1.png}{76.8}{0}
\qfigdef{p1_chem_s2_q11_A}{TNAS - Test 04 (18 Jun 2023)-img/p1_chem_s2_q11_A.png}{37.2}{0}
\qfigdef{p1_chem_s2_q11_B}{TNAS - Test 04 (18 Jun 2023)-img/p1_chem_s2_q11_B.png}{45.1}{0}
\qfigdef{p1_chem_s2_q12_1}{TNAS - Test 04 (18 Jun 2023)-img/p1_chem_s2_q12_1.png}{166.3}{0}
\qfigdef{p1_chem_s2_q12_4}{TNAS - Test 04 (18 Jun 2023)-img/p1_chem_s2_q12_4.png}{50.0}{0}
\qfigdef{p1_chem_s2_q12_2}{TNAS - Test 04 (18 Jun 2023)-img/p1_chem_s2_q12_2.png}{28.9}{0}
\qfigdef{p1_chem_s2_q12_3}{TNAS - Test 04 (18 Jun 2023)-img/p1_chem_s2_q12_3.png}{38.9}{0}
\qfigdef{p1_chem_s3_q03_1}{TNAS - Test 04 (18 Jun 2023)-img/p1_chem_s3_q03_1.png}{217.0}{0}
\qfigdef{p1_chem_s3_q05_1}{TNAS - Test 04 (18 Jun 2023)-img/p1_chem_s3_q05_1.png}{24.0}{0}
\qfigdef{p2_phy_s1_q01_1}{TNAS - Test 04 (18 Jun 2023)-img/p2_phy_s1_q01_1.png}{105.8}{0}
\qfigdef{p2_phy_s1_q02_1}{TNAS - Test 04 (18 Jun 2023)-img/p2_phy_s1_q02_1.png}{108.3}{0}
\qfigdef{p2_phy_s1_q04_1}{TNAS - Test 04 (18 Jun 2023)-img/p2_phy_s1_q04_1.png}{108.2}{0}
\qfigdef{p2_phy_s1_q06_1}{TNAS - Test 04 (18 Jun 2023)-img/p2_phy_s1_q06_1.png}{86.2}{0}
\qfigdef{p2_phy_s1_q07_1}{TNAS - Test 04 (18 Jun 2023)-img/p2_phy_s1_q07_1.png}{89.6}{0}
\qfigdef{p2_phy_s1_q08_1}{TNAS - Test 04 (18 Jun 2023)-img/p2_phy_s1_q08_1.png}{130.3}{0}
\qfigdef{p2_phy_s2_q09_1}{TNAS - Test 04 (18 Jun 2023)-img/p2_phy_s2_q09_1.png}{148.2}{0}
\qfigdef{p2_phy_s2_q10_1}{TNAS - Test 04 (18 Jun 2023)-img/p2_phy_s2_q10_1.png}{183.5}{0}
\qfigdef{p2_phy_s2_q11_1}{TNAS - Test 04 (18 Jun 2023)-img/p2_phy_s2_q11_1.png}{82.0}{0}
\qfigdef{p2_phy_s2_q12_1}{TNAS - Test 04 (18 Jun 2023)-img/p2_phy_s2_q12_1.png}{82.0}{0}
\qfigdef{p2_phy_s3_q01_1}{TNAS - Test 04 (18 Jun 2023)-img/p2_phy_s3_q01_1.png}{81.8}{0}
\qfigdef{p2_phy_s3_q03_1}{TNAS - Test 04 (18 Jun 2023)-img/p2_phy_s3_q03_1.png}{124.7}{0}
\qfigdef{p2_phy_s3_q06_1}{TNAS - Test 04 (18 Jun 2023)-img/p2_phy_s3_q06_1.png}{174.1}{0}
\qfigdef{p2_chem_s1_q07_1}{TNAS - Test 04 (18 Jun 2023)-img/p2_chem_s1_q07_1.png}{126.8}{0}
\qfigdef{p2_chem_s1_q07_2}{TNAS - Test 04 (18 Jun 2023)-img/p2_chem_s1_q07_2.png}{108.2}{0}
\qfigdef{p2_chem_s1_q07_3}{TNAS - Test 04 (18 Jun 2023)-img/p2_chem_s1_q07_3.png}{108.3}{0}
\qfigdef{p2_chem_s1_q07_4}{TNAS - Test 04 (18 Jun 2023)-img/p2_chem_s1_q07_4.png}{108.3}{0}
\qfigdef{p2_chem_s1_q07_5}{TNAS - Test 04 (18 Jun 2023)-img/p2_chem_s1_q07_5.png}{108.5}{0}
\qfigdef{p2_chem_s1_q08_A}{TNAS - Test 04 (18 Jun 2023)-img/p2_chem_s1_q08_A.png}{82.0}{0}
\qfigdef{p2_chem_s1_q08_2}{TNAS - Test 04 (18 Jun 2023)-img/p2_chem_s1_q08_2.png}{82.0}{0}
\qfigdef{p2_chem_s1_q08_B}{TNAS - Test 04 (18 Jun 2023)-img/p2_chem_s1_q08_B.png}{71.7}{0}
\qfigdef{p2_chem_s1_q08_4}{TNAS - Test 04 (18 Jun 2023)-img/p2_chem_s1_q08_4.png}{44.9}{0}
\qfigdef{p2_chem_s1_q08_C}{TNAS - Test 04 (18 Jun 2023)-img/p2_chem_s1_q08_C.png}{87.0}{0}
\qfigdef{p2_chem_s1_q08_6}{TNAS - Test 04 (18 Jun 2023)-img/p2_chem_s1_q08_6.png}{82.1}{0}
\qfigdef{p2_chem_s2_q11_1}{TNAS - Test 04 (18 Jun 2023)-img/p2_chem_s2_q11_1.png}{119.3}{0}
\qfigdef{p2_chem_s2_q11_2}{TNAS - Test 04 (18 Jun 2023)-img/p2_chem_s2_q11_2.png}{118.7}{0}
\qfigdef{p2_chem_s3_q04_1}{TNAS - Test 04 (18 Jun 2023)-img/p2_chem_s3_q04_1.png}{77.0}{0}
\qfigdef{p2_chem_s3_q05_1}{TNAS - Test 04 (18 Jun 2023)-img/p2_chem_s3_q05_1.png}{175.3}{0}
\qfigdef{p2_chem_s3_q05_2}{TNAS - Test 04 (18 Jun 2023)-img/p2_chem_s3_q05_2.png}{179.5}{0}
\qfigdef{p2_chem_s3_q05_3}{TNAS - Test 04 (18 Jun 2023)-img/p2_chem_s3_q05_3.png}{280.4}{1}
\qfigdef{p2_math_s1_q08_1}{TNAS - Test 04 (18 Jun 2023)-img/p2_math_s1_q08_1.png}{79.9}{0}
\qfigdef{p2_math_s3_q06_1}{TNAS - Test 04 (18 Jun 2023)-img/p2_math_s3_q06_1.png}{128.8}{0}
\raggedcolumns
\makeatletter\@beginparpenalty=10000\makeatother
\newcolumntype{L}[1]{>{\raggedright\arraybackslash}p{#1}}
\begin{document}
\begin{center}{\Large\bfseries TNAS --- Test 04 (18 Jun 2023)}\end{center}
\begin{multicols}{2}
\raggedcolumns
\papersec{1}
\partsec{1}{PHYSICS}
\sect{SECTION-I(i)}
\qn{1} A light uniform rod of young's modulus $Y$, cross sectional area $A$, coefficient of linear expansion $\alpha$ and length $\ell_0$ is rigidly connected to support at one end and the other end of the rod is connected to the spring shown in the figure. The temperature of the rod increased by $\Delta\theta$ with supports remaining fixed. Initially the spring is in natural length position. Spring force on the rod acts uniformly over the cross section during elongation of the rod. Find the net elongation of the rod ? (Assume thermal strain to be small)
\qfig{p1_phy_s1_q01_1}
\begin{qoptions}
\opt{A}{$\dfrac{k\ell_0\alpha\Delta\theta}{k+\frac{YA}{\ell_0}}$}
\opt{B}{$\dfrac{YA\alpha\Delta\theta}{k+\frac{YA}{\ell_0}}$}
\opt{C}{$YA\alpha\Delta\theta$}
\opt{D}{None of these}
\end{qoptions}
\qn{2} The ends of a rod of length $\ell$ are maintained at $\theta_1$ and $\theta_2$ ($<\theta_1$). The thermal conductivity of rod at a point is given as $k=\dfrac{\alpha}{\theta}$ where $\theta$ = temperature at that point of the rod $\theta$ as a function of $x$ is.
\qfig{p1_phy_s1_q02_1}
\begin{qoptions}
\opt{A}{$\dfrac{x}{\ell}\dfrac{\theta_1^2}{\theta_2}$}
\opt{B}{$\dfrac{\theta_1 x}{\ell}e^{\theta_1/\theta_2}$}
\opt{C}{$\left(\dfrac{x}{\ell}\right)^2\dfrac{\theta_1^2}{\theta_2^2}$}
\opt{D}{$\theta_1\left(\dfrac{\theta_2}{\theta_1}\right)^{x/\ell}$}
\end{qoptions}
\qn{3} A hollow cube of side $L$ has two parallel sides conducting and remaining four sides are non conducting. It is filled with a liquid of dielectric constant $k$. A small hole of area $A$ at the bottom of the cube is opened at $t=0$ and liquid starts leaking through it. The two conducting sides are connected by a battery of emf $\varepsilon$. [$L=9$ m, $A=\dfrac{18}{\sqrt{20}}\,\mathrm{m}^2$, $k=10$, $\varepsilon=81$ v]. Find the capacitance of capacitor as a function of time ?
\qfig{p1_phy_s1_q03_1}
\begin{qoptions}
\opt{A}{$\dfrac{\varepsilon_0}{9}\left(t^2-54t+64\right)$}
\opt{B}{$\dfrac{\varepsilon_0}{9}\left(t^2-54t+810\right)$}
\opt{C}{$\dfrac{\varepsilon_0}{9}\left(t^2+54t+810\right)$}
\opt{D}{$\dfrac{\varepsilon_0}{9}\left(t^2+54t+64\right)$}
\end{qoptions}
\qn{4} An alternating voltage $E=6\sin(20t)+8\cos(20t)$ is applied to a series resonate circuit as shown. Mark the \textbf{INCORRECT} statement.
\qfig{p1_phy_s1_q04_1}
\begin{qoptions}
\opt{A}{The current in the circuit is 2A.}
\opt{B}{Power dissipated in the circuit is 10 W.}
\opt{C}{The quality factor of the circuit is 0.8.}
\opt{D}{The capacitance C is 12.5 mF.}
\end{qoptions}
\sect{SECTION-I(ii)}
\qn{5} A rectangle of perfectly conducting wire having sides $a$ and $b(<a)$, mass $M$, self inductance $L$, moves with an initially velocity $v_0$ in its plane directed along its longer side, from a region of zero magnetic field into a region with a field $B_0$, which is uniform and perpendicular to the plane of rectangle. At $t=0$, rectangle just starts to enter magnetic field. Choose the correct statement(s):-
\begin{qoptions}
\opt{A}{Velocity of rectangle as a function of time is given by $v=v_0\cos\omega t$, where $\omega=\dfrac{B_0a}{\sqrt{mL}}$}
\opt{B}{Velocity of rectangle as a function of time is given by $v=v_0(1+\sin\omega t)$, where $\omega=\dfrac{B_0a}{\sqrt{mL}}$.}
\opt{C}{Displacement of rectangle at any time $t$ is $S=\dfrac{v_0\sqrt{mL}}{B_0a}\sin\omega t$, where $\omega=\dfrac{B_0a}{\sqrt{mL}}$.}
\opt{D}{Displacement of rectangle at any time $t$ is $S=\dfrac{v_0\sqrt{mL}}{B_0a}(1+\sin\omega t)$ $\omega=\dfrac{B_0a}{\sqrt{mL}}$}
\end{qoptions}
\qn{6} A particle of mass $m$ and charge $q$ starts moving from rest from the origin in crossed electric and magnetic field. The electric field has a constant magnitude $E$ and points towards positive $y$ direction. The magnetic field points towards positive $z$ direction but the magnitude of magnetic field changes with $y$ as $\left|\vec B\right|=\alpha\sqrt y$, where $\alpha$ is constant. The particle comes to rest for the first time at $t=T$.
\begin{qoptions}
\opt{A}{The speed of the particle at the moment when its $y$-coordinate is maximum is $\sqrt{\dfrac{3E}{\alpha}\sqrt{\dfrac{2qE}{m}}}$.}
\opt{B}{The radius of curvature of particle when its $y$-coordinate is maximum is $\dfrac{3}{\alpha}\sqrt{\dfrac{Em}{2q}}$.}
\opt{C}{The particle moves in a semicircle before coming to rest at $t=T$.}
\opt{D}{The displacement of particle at $t=\dfrac32T$ is $\dfrac{3}{\alpha}\sqrt{\dfrac{5Em}{2q}}$.}
\end{qoptions}
\qn{7} The circuit shown in figure (a) can be reduced into an equivalent current source ($I_t$) and parallel resistance ($R_t$) as shown in figure (b).
\qfig{p1_phy_s2_q07_1}
\qfig{p1_phy_s2_q07_2}
\begin{qoptions}
\opt{A}{The value of $I_t$ is $\dfrac{I_S}{9}$.}
\opt{B}{The value of $I_t$ is $\dfrac{2I_S}{9}$.}
\opt{C}{The value of $R_t$ is $9R$}
\opt{D}{The value of $R_t$ is $3R$}
\end{qoptions}
\qn{8} The 2.5 kg spring supported cylinder is set into vertical free oscillation and is observed to have a time period of 0.25 sec. The system is then immersed in a complete oil bath as shown in figure (b) and cylinder is again set in vertical oscillation. Due to damping it is observed that the amplitude of oscillation decreases by a factor of 4 after every oscillation.
\qfig{p1_phy_s2_q08_1}
\begin{qoptions}
\opt{A}{The equivalent spring constant of system of springs is 400 N/m.}
\opt{B}{The damping constant due to resistive force of oil (approximately) is $\dfrac{40\pi}{3}$ kg/s.}
\opt{C}{The time period of oscillation when system is placed in oil is (approximately) is 0.26 sec.}
\opt{D}{The energy of oscillation also decreases by a factor of 4 after every oscillation.}
\end{qoptions}
\qn{9} For the circuit consisting of infinite resistors in both directions as shown in the diagram, each resistance has value $R$.
\qfig{p1_phy_s2_q09_1}
\begin{qoptions}
\opt{A}{The equivalent resistance between A and B is $\dfrac{6-\sqrt3}{6}R$.}
\opt{B}{The equivalent resistance between A and B is $\dfrac{\sqrt2-1}{\sqrt2}R$.}
\opt{C}{The equivalent resistance between A and C is $\dfrac{R}{\sqrt3}$.}
\opt{D}{The equivalent resistance between B and C is $\sqrt3R$}
\end{qoptions}
\qn{10} The reading of voltmeter is 60V and reading of ammeter $A_1$ is 2A. The power loss across $R_3$ is 250 W. Then
\qfig{p1_phy_s2_q10_1}
\begin{qoptions}
\opt{A}{The values of resistance $R_2$ is 20 $\Omega$.}
\opt{B}{The value of resistance $R_3$ is $10\Omega$.}
\opt{C}{The reading of ammeter $A_2$ is 5A.}
\opt{D}{The EMF of battery is 100 V.}
\end{qoptions}
\qn{11} A heat exchanger is a device which transfers heat between a hot fluid and a cold fluid. The heat exchanger shown below consists of two rectangular tubes of length $\ell$, width $w$ and height $h$. The tubes are separated by a metal wall of thickness $d$ and thermal conductivity $k$. Originally hot fluid flows through the lower tube at a speed $v$ from right to left, and originally cold fluid flows through the upper tube in the opposite direction (left to right) at the same speed. The heat capacity per unit volume of both fluids is $c$.\par
The hot fluid enters the heat exchanger at a higher temperature than the cold fluid; the difference between the temperatures of the entering fluids is $\Delta T_i$. When the fluids exit the heat exchanger the difference has been reduced to $\Delta T_f$.
\qfig{p1_phy_s2_q11_1}
Assume that the temperature in each pipe changes linearly with the length of the tube, and consider transfer of heat only due to conduction in the metal and due to the bulk movement of fluid. Under the assumptions in this problem, while the temperature of each fluid varies along the length of the exchanger, the temperature difference across the wall is the same everywhere.
\begin{qoptions}
\opt{A}{The amount of energy transferred from one liquid to another per unit time is $\dfrac{chwv\Delta T_i}{\left(1+\frac{chvd}{\ell k}\right)}$.}
\opt{B}{The change in temperature of a liquid as it crosses the length $\ell$ is $\dfrac{\Delta T_i}{\left(1+\frac{chvd}{\ell k}\right)}$.}
\opt{C}{The increase in temperature of one liquid is equal to decrease in temperature of other liquid.}
\opt{D}{The value of $\Delta T_f$ is $\left[1-\dfrac{2}{1+\frac{chvd}{\ell k}}\right]\Delta T_i$}
\end{qoptions}
\qn{12} A uniform heavy vertical rope hanging under its own weight is being pulled will an acceleration $a=g$ as shown in the figure. At $t=0$, velocity of rope is zero and length of vertical rope is $\ell$. As rope begins to move, a small transverse fulse is generated at the lowest point.
\qfig{p1_phy_s2_q12_1}
\begin{qoptions}
\opt{A}{The speed of pulse wrt rope as a function of $x$ is $\sqrt{2gx}$.}
\opt{B}{The speed of pulse wrt rope as a function of $x$ is $\sqrt{gx}$.}
\opt{C}{The time in which the pulse reaches the pulley is $\sqrt{\dfrac{2\ell}{g}}$.}
\opt{D}{The time in which the pulse reaches the pulley is $\sqrt{\dfrac{\ell}{g}}$.}
\end{qoptions}
\sect{SECTION-II}
\qn{1} Water of mass $m$ and specific heat capacity $S$ is kept at $60^\circ$ C inside a thick spherical shell of inner and outer radii $a$ and $2a$ as shown. Thermal conductivity of the shell varies as $k=\dfrac{C}{r^3}$, where $C$ is a constant and $r$ is the distance from the centre. The surrounding temperature is $\theta_0=30^\circ$C. If the time required to decrease the temperature of water from $60^\circ$C to $40^\circ$C is $t=n\left(\dfrac{mSa^3\ell n3}{8\pi C}\right)$. Find the value of $n$.
\qfig{p1_phy_s3_q01_1}
\qn{2} A sphere made of polymer is floating on the surface of water which is at a temperature of $20^\circ$C and 1\% of the volume of the sphere is above the water. What final temperature (in $^\circ$C) must the water be heated in order to submerge the polymer sphere completely. Take coefficient of linear expansion of the polymer $2.30\times10^{-5}$ /$^\circ$C and coefficient of volume expansion of water is $3.19\times10^{-4}$ /$^\circ$C.
\qn{3} A galvanometer of resistance $98\Omega$ is connected in a circuit as shown in figure. It has 40 divisions on both sides of zero and it is showing full scale deflection in the given situation. The figure of merit of galvanometer is (in $\mu$A / division)
\qfig{p1_phy_s3_q03_1}
\qn{4} Metal rod with a mass $m=10$ g, and length $L=0.2$ m is suspended by two light wires length $\ell=1$ cm in a magnetic field induction $B$ which is directed vertically downwards (Figure). A capacitor of capacitance $C=100\ \mu$F charged to a voltage of 100 V is connected as shown. The maximum deflection (in degree) of the rod from the initial position after the switch is closed is $37^\circ$. Resistance of wire and rod is not taken into account. What is the magnitude of magnetic field in Tesla ? (Given $g=10$ m/s$^{-2}$)
\qfig{p1_phy_s3_q04_1}
\qn{5} Two pieces of copper wires are soldered together, such that the two pieces form semicircles and together the wires form a circle of radius $r=4$cm. The diameters of the wires are $d_1=3$ mm and $d_2=1.5$ mm. To one of the solder points of the closed circle (A) and to the midpoint of the semicircle made of thinner wire (C) very long straight wires are connected (one to each point). The magnetic induction at the centre of the circular wire, when the current in the straight wires is $I=25$ A, is $\alpha\times10^{-5}$ T. The value of $\alpha$ is
\qfig{p1_phy_s3_q05_1}
\qn{6} Switch K in the circuit shown in the figure has been closed for a long time. The magnitude of the induced electromotive force (in V) in the coil immediately after the switch is opened is
\qfig{p1_phy_s3_q06_1}
\partsec{2}{CHEMISTRY}
\sect{SECTION-I(i)}
\qn{1} \qfig{p1_chem_s1_q01_1}
Find the total number of moles of A formed.
\begin{qoptions}
\opt{A}{6}
\opt{B}{2}
\opt{C}{4}
\opt{D}{can not be predicted}
\end{qoptions}
\qn{2} The total number of cis $\mathrm{N-\widehat{Co}-Cl}$ bond angle in cis-$\mathrm{[Co(en)_2Cl_2]^{+1}}$ is.
\begin{qoptions}
\opt{A}{8}
\opt{B}{6}
\opt{C}{4}
\opt{D}{None of these}
\end{qoptions}
\qn{3} \qfig{p1_chem_s1_q03_1}
\qfig{p1_chem_s1_q03_2}
Relation between (A) and (B) is
\begin{qoptions}
\opt{A}{Geometrical isomer}
\opt{B}{Positional isomer}
\opt{C}{Enantiomer}
\opt{D}{Chain isomer}
\end{qoptions}
\qn{4} \qfig{p1_chem_s1_q04_1}
Identify correct statement for ``Q'' product?
\begin{qoptions}
\opt{A}{``Q'' product is benzaldehyde}
\opt{B}{``Q'' product is acetophenone}
\opt{C}{``Q'' product is benzyl alcohol}
\opt{D}{``Q'' product is benzoic acid.}
\end{qoptions}
\sect{SECTION-I(ii)}
\qn{5} $\mathrm{2A(g)\longrightarrow 4B(g)+C(g)}$\par
At $t=0$ pressure of A is 2 atm, at $t=10$ sec. pressure of B is 2 atm. At $t=30$ sec. pressure of C is 0.75 atm. Select correct option/options :
\begin{qoptions}
\opt{A}{At $t=70$ sec pressure of B will be 3.5 atm.}
\opt{B}{Reaction may be an elementary reaction}
\opt{C}{At $t=70$ sec rate of formation of B will be $\dfrac{1}{80}$ atm/sec.}
\opt{D}{Reaction will complete in definite time.}
\end{qoptions}
\qn{6} A compound $\mathrm{AB_2}$ crystallizes similarly as $\mathrm{CaF_2}$. $r_{A^{+2}}=$ `x' $r_{B^-}=$ `y' where $r_{A^{+2}}$ and $r_{B^-}$ are respective radii. Select the correct statement(s) for the above crystal :
\begin{qoptions}
\opt{A}{Minimum distance between two $\mathrm{B^-}$ ions will be $\dfrac{2}{\sqrt3}(x+y)$}
\opt{B}{Two $\mathrm{B^-}$ ions may be present at a distance of $2(x+y)$}
\opt{C}{One $\mathrm{A^{+2}}$ ion is surrounded by eight $\mathrm{B^-}$ ions in lattice}
\opt{D}{It can show anion Frenkel defect}
\end{qoptions}
\qn{7} Choose correct statement(s) :
\begin{qoptions}
\opt{A}{All activated complex get converted into products.}
\opt{B}{Half life of `A', would be independent of initial concentration of `B' in an elementary reaction : $\mathrm{A(g)+2B(g)\longrightarrow}$ Products which starts with 1 mol `A' and 1000 moles of B in 1 litre vessel.}
\opt{C}{In I order reaction : $\mathrm{A\xrightarrow{K_1}B\xrightarrow{K_2}C}$ rate of formation of C will be maximum at time $\dfrac{\ln(K_2/K_1)}{(K_2-K_1)}$}
\opt{D}{A parallel first order reaction \qfigi{p1_chem_s2_q07_1} $K_1:K_2=\dfrac21$. Mole ratio of `B' \& `C' would be $3:2$, when 50\% of `A' has reacted.}
\end{qoptions}
\qn{8} In the mixture of two gases `A' and `B' rate of effusion through a pin hole of gases `A' and `B' are in the ratio of 16: 3. If the initial ratio of their masses present in the mixture is 2 : 3. Then
\begin{qoptions}
\opt{A}{The ratio of their molar masses is 16 : 1}
\opt{B}{The ratio of their molar masses is 1 : 4}
\opt{C}{The ratio of their moles present inside the container is 1 : 24 initially}
\opt{D}{The ratio of their moles present inside the container is 8 : 3 initially}
\end{qoptions}
\qn{9} When $\mathrm{3F^-}$ is replaced by $\mathrm{3H_2O}$ molecule from Hexafluoridocobaltate (III) then which of the forming property changed.
\begin{qoptions}
\opt{A}{Hybridisation}
\opt{B}{$\Delta_0$}
\opt{C}{Magnetic moment}
\opt{D}{Number of $\mathrm{t_{2g}}$ electron}
\end{qoptions}
\qn{10} A $\mathrm{Pt^{+2}}$ complex having EAN value 84 and having two type of achiral monodentate ligand of $\mathrm{Pt^{2+}}$ complex. If two ligands are replaced from trans isomer by one butane 2, 3-diamine then which of the following property changed.
\begin{qoptions}
\opt{A}{Number of geometrical isomer only}
\opt{B}{Number of optical isomer only}
\opt{C}{(A) \& (B) both}
\opt{D}{None of these}
\end{qoptions}
\qn{11} \qfig{p1_chem_s2_q11_1}
Which of the following product can be formed in the given reaction?
\begin{qoptions}
\opt{A}{\qfigi{p1_chem_s2_q11_A}}
\opt{B}{\qfigi{p1_chem_s2_q11_B}}
\opt{C}{$\mathrm{Et_3N}$}
\opt{D}{$\mathrm{CH_2=CH_2}$}
\end{qoptions}
\qn{12} \qfig{p1_chem_s2_q12_1}
\begin{qoptions}
\opt{A}{D is $\mathrm{CH_3{-}CH{=}CH_2}$}
\opt{B}{E is \qfigi{p1_chem_s2_q12_4}}
\opt{C}{F is \qfigi{p1_chem_s2_q12_2}}
\opt{D}{B is \qfigi{p1_chem_s2_q12_3}}
\end{qoptions}
\sect{SECTION-II}
\qn{1} A sample of camphor used in a method for determining of molar mass had melting point of 450 K. The melting point of a solution containing 0.5 gm camphor and 0.04 gm of an unknown non-volatile non-electrolyte compound was 430 K. The molar mass of camphor is 152 and its $\Delta H_{\mathrm{fusion}}=1.52$ Kcal mol$^{-1}$. The molar mass of unknown compound in gm/mole is \dots\dots\dots\\
(R = 2 cal mol$^{-1}$ K$^{-1}$).
\qn{2} Find the inverse of critical compressibility factor $\left(\dfrac{1}{Z_C}\right)$ for the following equation of state.\par
$PV=RT-\dfrac{A}{V}+\dfrac{2B}{V^2}$\par
where A and B are constant, $P$ = pressure and $V$ = molar volume.
\qn{3} When NO reacts with aqueous solution of $\mathrm{FeSO_4}$ then brown colour complex is formed then find the number of unpaired electron in non axial set of `d' orbital of metal.
\qn{4} \qfig{p1_chem_s3_q03_1}
\qn{5} In the given reaction the value of x is?
\qfig{p1_chem_s3_q05_1}
{\small Monochlorination: $\mathrm{Cl_2/h\nu}$ $\to$ (A) $\xrightarrow{\mathrm{Na/ether}}$ (B) $\xrightarrow[\text{monochlorination}]{\mathrm{Cl_2/h\nu}}$ total product $\xrightarrow{\text{fractional distillation}}$ $m$ (number of fraction)}\par
what is the value of `m'?
\qn{6} No. of enantiomeric pair of $\mathrm{[Ma_2bcde]}$ complex = P\\
No. of geometrical isomers for $\mathrm{[PdCl_3Br(SCN)(NO_2)]^{2-}}$ = Q\\
No. of unpaired electron in non axial d-orbital of $\mathrm{[Co(H_2O)_6]^{3+}}$ = R\\
Sum of P + Q + R is :
\partsec{3}{MATHEMATICS}
\sect{SECTION-I(i)}
\qn{1} If $y=\dfrac{ax+b}{cx+d}$, then value of $\dfrac{y_2^2}{y_1y_3}$, wherever exists, is (where $y_n$ denotes $n^{\mathrm{th}}$ order derivative of differentiable function $y$ and $a,b,c,d$ are constants) -
\begin{qoptions}
\opt{A}{$\dfrac32$}
\opt{B}{$\dfrac23$}
\opt{C}{$\dfrac45$}
\opt{D}{$\dfrac54$}
\end{qoptions}
\qn{2} If $f(x)=\begin{vmatrix}(x-1)^2&(x-2)^3&(x+1)\\(x-1)^3&x^2&(x+1)^2\\x+2&x+1&x^5\end{vmatrix}$ then coefficient of $x$ in $f(x)$ is -
\begin{qoptions}
\opt{A}{$-14$}
\opt{B}{$-16$}
\opt{C}{0}
\opt{D}{25}
\end{qoptions}
\qn{3} {\footnotesize$\displaystyle\int\left(2e^{\tan^{-1}x+x^2+\ln x}+\dfrac{(2x^3+2x+1)}{\left(x+\frac1x\right)}e^{\tan^{-1}x+x^2+\ln x}\right)dx$} is equal to
\begin{qoptions}
\opt{A}{$\dfrac{x^3}{3}e^{\tan^{-1}x+x^2+\ln x}+C$}
\opt{B}{$\dfrac{x^2}{2}e^{\tan^{-1}x+x^2+\ln x}+C$}
\opt{C}{$xe^{\tan^{-1}x+x^3+\ln x}+C$}
\opt{D}{$x^2e^{\tan^{-1}x+x^2}+C$ (where $C$ is constant of integration)}
\end{qoptions}
\qn{4} The function $f(x)=\alpha x^2-\beta x-4x^3+\gamma$, where $\alpha,\beta,\gamma\in R$, has local minima at $P(\log_2a,f(\log_2a))$ \& local maxima at $Q(\log_2a^2,f(\log_2a^2))$. If the graph of $f(x)$ changes concavity about the point $R\left(\dfrac34,f\left(\dfrac34\right)\right)$, then which of the following conic section can have eccentricity `$a$' -
\begin{qoptions}
\opt{A}{circle}
\opt{B}{parabola}
\opt{C}{ellipse}
\opt{D}{hyperbola}
\end{qoptions}
\sect{SECTION-I(ii)}
\qn{5} $ABC$ is a triangle in which $\angle C=2\angle B$. $D$ is a point on $BC$ such that $AD$ bisects $\angle BAC$ and $AC=BD$. Then which of the following is/are true ?
\begin{qoptions}
\opt{A}{$\cos B+\cos C=\dfrac{\sqrt5}{2}$}
\opt{B}{$\cos B-\cos C=\dfrac12$}
\opt{C}{$\dfrac{\sin A}{\sin C}=1$}
\opt{D}{$\dfrac{\sin A}{\sin C}=\dfrac{\sqrt5+1}{4}$}
\end{qoptions}
\qn{6} Which of the following statements is/are always true ?
\begin{qoptions}
\opt{A}{number of roots of equation $\displaystyle\sum_{r=1}^{2017}\dfrac{r}{x-r}=0$ is 2016}
\opt{B}{number of roots of equation $\displaystyle\sum_{r=1}^{2017}\dfrac{r}{x-r}=0$ is 2017}
\opt{C}{if $f(x)$ is continuous in $[0,2]$ such that $f(0)=f(2)$, then there exists atleast one $x\in[0,1]$ such that $f(x)=f(x+1)$.}
\opt{D}{if $f(x)$ is continuous in $[0,2]$ such that $f(0)=f(2)$, then there exists atleast one $x\in(0,1)$ such that $f(x)=f(x+1)$.}
\end{qoptions}
\qn{7} If $a+b+c=4$ and $a^2+b^2+c^2+3(ab+bc+ac)=21$ where $a,b,c\in R$, then which of the following is/are true ?
\begin{qoptions}
\opt{A}{minimum value of $(a+b)(b+c)(a+c)$ is 18}
\opt{B}{minimum value of $(a+b)(b+c)(a+c)$ is 24}
\opt{C}{maximum value of $(a+b)(b+c)(a+c)$ is $\dfrac{490}{27}$}
\opt{D}{maximum value of $(a+b)(b+c)(a+c)$ is $\dfrac{471}{27}$}
\end{qoptions}
\qn{8} Let $f(x)=\begin{cases}x[x]&-4\leqslant x<-1\\1-x&-1\leqslant x<1\\x^2-3x+2&1\leqslant x<2\\(x-2)[x]&2\leqslant x\leqslant4\end{cases}$, then (where $[.]$ denotes greatest integer function)
\begin{qoptions}
\opt{A}{number of points where $f(x)$ is non-differentiable is 6}
\opt{B}{sum of global maximum value and global minimum value of $f(x)$ is $\dfrac{63}{4}$}
\opt{C}{number of points of extrema of $f(x)$ is 3 points}
\opt{D}{number of solutions of $f(x)=4$ is 3}
\end{qoptions}
\qn{9} Let $F(n,x)=\displaystyle\sum_{r=1}^{n}\tan^{-1}\left(\dfrac{2x}{4+r(r+1)x^2}\right)$, where $x>0$ and $n\in N$ and $\dfrac{d}{dx}F(n,x)=G(n,x)$, then
\begin{qoptions}
\opt{A}{$F(8,2)=\tan^{-1}\left(\dfrac45\right)$}
\opt{B}{$\displaystyle\lim_{n\to\infty}F\left(n,2\sqrt3\right)=\dfrac\pi6$}
\opt{C}{If $G\left(n,\dfrac13\right)<0$, then $n\geq36$}
\opt{D}{$\displaystyle\lim_{n\to\infty}G(n,1)=-\dfrac25$}
\end{qoptions}
\qn{10} Let $f(x)$ be a non constant twice derivable function defined on $R$ such that $f(2+x)=f(2-x)$ and $f'\left(\dfrac12\right)=0=f'(1)$, then minimum number of roots of the equation $f''(x)=0$ in $(0,4)$, is
\begin{qoptions}
\opt{A}{4}
\opt{B}{3}
\opt{C}{2}
\opt{D}{5}
\end{qoptions}
\qn{11} Consider the following system of equations\\
$\alpha x+y+z=m$\\
$x+\alpha y+z=n$\\
$x+y+\alpha z=p$\\
then which of the following statement is/are correct -
\begin{qoptions}
\opt{A}{The given system of equations has no solution if $\alpha=-2$ and $m+n+p\neq0$}
\opt{B}{The given system of equations has no solution if $\alpha=1$ and $m\neq n$ or $n\neq p$ or $p\neq m$}
\opt{C}{The given system of equations has infinite solution if $\alpha=-2$ and $m+n+p=0$}
\opt{D}{The given system of equations has unique solution if $\alpha=1$ or $-2$}
\end{qoptions}
\qn{12} If the sides $a,b,c$ of a triangle are 3, 5 and 7 respectively then which of the following is/are true ?
\begin{qoptions}
\opt{A}{$\tan\left(\dfrac{A-B}{2}\right)=\dfrac{-1}{4\sqrt3}$}
\opt{B}{largest angle is $120^\circ$}
\opt{C}{length of median through vertex $A$ is $\dfrac{\sqrt{139}}{4}$}
\opt{D}{length of angle bisector through $C$ is $\dfrac{15}{8}$}
\end{qoptions}
\sect{SECTION-II}
\qn{1} If $\alpha,\beta,\gamma$ are the roots of $4x^3-16x^2+19x-6=0$, then the value of the determinant $\begin{vmatrix}\alpha-\beta-\gamma&2\alpha&2\alpha\\2\beta&\beta-\gamma-\alpha&2\beta\\2\gamma&2\gamma&\gamma-\alpha-\beta\end{vmatrix}$ is
\qn{2} A curve is represented by $x=\sec^2t$ and $y=\cot t$, where $t$ is the parameter. If the tangent at point $P$ on curve where $t=\dfrac\pi4$ meets the curve again at $Q$ then $|PQ|=\dfrac{p\sqrt q}{2}$, where $p,q$ are coprime. Find $p+q$
\qn{3} Let $f$ be a real valued function such that $f\left(\dfrac{2x-3}{x-2}\right)=5x-2$, $x\neq2$, then value of $\dfrac54f^{-1}(13)$ is
\qn{4} In a triangle $ABC$, medians $AD$ and $CE$ are drawn. If $AD=5$, $\angle DAC=\dfrac\pi8$ and $\angle ACE=\dfrac\pi4$ then the area of the triangle $ABC$ is equal to
\qn{5} Let $p_1$ and $p_2$ be two values of $p$ for which $f(x)=\begin{cases}\dfrac{x(\ln(1+x)+\ln(1-x))}{\sec x-\cos x};&x\in(-1,0)\\(p^2-3p-1)\sin x+x^2;&x\in[0,\infty)\end{cases}$ is differentiable at $x=0$ find $p_1^2+p_2^2$
\qn{6} Let $P(x)$ be a polynomial of degree 4 with $P(1)=1$, $P'(1)=0$, $P''(1)=2$, $P'''(1)=-12$ and $P''''(1)=24$ then value of $\dfrac{P(2)}{P(3)}$ is equal to
\papersec{2}
\partsec{1}{PHYSICS}
\sect{SECTION-I(i)}
\qn{1} In the given figure `P' is a point charge of charge `q' placed inside on arbitrary shaped cavity of an uncharged conducting sphere of radius $R$ at a distance $R/2$ from centre `C', `S' is another point charge of charge ``2q'' placed at distance $2R$ from centre. PCS is a straight line. Choose this correct option(s) :-
\qfig{p2_phy_s1_q01_1}
\begin{qoptions}
\opt{A}{P applies force on outer surface.}
\opt{B}{S applied force on inner surface.}
\opt{C}{Net electric force on P due to charge on outer surface is $\dfrac{2}{25}\dfrac{q^2}{\pi\epsilon_0R^2}$.}
\opt{D}{Net electric force on S due to ``sphere + P'' is rightward.}
\end{qoptions}
\qn{2} An infinite wire, placed along $z$-axis has current $I_1$ in positive $z$-direction. A conducting rod $PQ$ placed in $xy$ plane parallel to $y$-axis has current $I_2$ in positive $y$-direction. The ends of the rod subtend angles $37^\circ$ and $53^\circ$ at the origin with positive $x$-axis as shown in figure. The rod is at distance $\ell$ from origin. Then choose the correct options :-
\qfig{p2_phy_s1_q02_1}
\begin{qoptions}
\opt{A}{The direction of net force acting on the wire $PQ$ is along the negative $z$-axis.}
\opt{B}{The direction of net force acting on the wire $PQ$ is along the positive $z$-axis.}
\opt{C}{The magnitude of net force acting on the rod $PQ$ is $\dfrac{\mu_0I_1I_2}{4\pi}\ell n\left(\dfrac{16}{9}\right)$.}
\opt{D}{The magnitude of net force acting on the rod $PQ$ is $\dfrac{\mu_0I_1I_2}{2\pi}\ell n\left(\dfrac{16}{9}\right)$}
\end{qoptions}
\qn{3} A voltmeter of resistance $R_1$ and an ammeter of resistance $R_2$ are connected in series across a battery of emf $V=450$ V having internal resistance $5\Omega$. When a resistance $12\Omega$ is connected in parallel to voltmeter, reading of ammeter increases four times while that of voltmeter reduces to one fourth. Then choose the correct options :-
\begin{qoptions}
\opt{A}{The resistor $R_1=180\ \Omega$}
\opt{B}{The resistor $R_2=40\Omega$.}
\opt{C}{Current read by ammeter is 3A when $12\ \Omega$ resistance is not connected.}
\opt{D}{Current read by ammeter is 3A when $12\ \Omega$ resistance is connected.}
\end{qoptions}
\qn{4} In the circuit shown in the figure, the key K is open and all capacitors are uncharged. At $t=0$, key is closed :-
\qfig{p2_phy_s1_q04_1}
\begin{qoptions}
\opt{A}{Reading of voltmeter is $\dfrac35U_0$.}
\opt{B}{Charge flown through the ammeter is $\dfrac35CU_0$.}
\opt{C}{Energy stored in all the capacitor is $\dfrac75CU_0^2$}
\opt{D}{None of the above}
\end{qoptions}
\qn{5} A thin hollow spherical shell of radius $R$ is filled with water. The temperature of atmosphere is $T_s$ which is below $0^\circ$C. Firstly all water comes down to $0^\circ$C and then starts freezing from the outer surface towards the centre. Assume that the thermal conductivity of shell is very high. Neglect change in density due to freezing. $\rho$ = density of ice, $k$ = thermal conductivity of ice, $L$ = latent heat of fusion.
\begin{qoptions}
\opt{A}{Time taken to whole water to freeze is $\dfrac{\rho LR^2}{6kT_S}$}
\opt{B}{Time taken to whole water to freeze is $\dfrac{2\rho LR^2}{kT_S}$.}
\opt{C}{Time taken to freeze first $\dfrac78$th of total mass is more than time taken for last $\dfrac M8$ mass to freeze.}
\opt{D}{Time taken to freeze first $\dfrac78$th of total mass is same as time taken for last $\dfrac M8$ mass to freeze.}
\end{qoptions}
\qn{6} The figure shows a circular loop made of a wire. The resistivity of the material varies as a function of $\theta$ such that $\rho=\rho_0\sin^2\theta$. The positions of jockey $J$ such that the magnetic field at the centre ($O$) due to the current in the loop is zero, will be
\qfig{p2_phy_s1_q06_1}
\begin{qoptions}
\opt{A}{$\theta=\dfrac\pi2$}
\opt{B}{$\theta=\pi$}
\opt{C}{$\theta=\dfrac{3\pi}{2}$}
\opt{D}{$\theta=\dfrac\pi4$}
\end{qoptions}
\qn{7} In the shown circuit initially, the switch is open and the circuit is in steady state then :-
\qfig{p2_phy_s1_q07_1}
\begin{qoptions}
\opt{A}{The charge on $3\mu$F capacitor in steady state before closing the switch is 20 $\mu$C.}
\opt{B}{The current through the battery in steady state before closing the switch is zero.}
\opt{C}{The charge on 2 $\mu$F in steady state after closing the switch is 8 $\mu$C.}
\opt{D}{The current through the battery in steady state after closing the switch is 2A.}
\end{qoptions}
\qn{8} An ideal battery of emf = 24 V is connected to the primary winding of an ideal transformer with a resistance $R_1=500\ \Omega$ and a switch. A resistance $R_2=300\ \Omega$ is connected to the secondary winding as shown the diagram. The primary coil has 800 turns and secondary coil has 1000 turns. Initially the switch K is open. Just after the switch is closed.
\qfig{p2_phy_s1_q08_1}
\begin{qoptions}
\opt{A}{The current induced in primary circuit (approx.) is 40 mA.}
\opt{B}{The current induced in primary circuit (approx.) is 35 mA.}
\opt{C}{The current induced in secondary circuit (approx.) is 28 mA.}
\opt{D}{The current induced in secondary circuit (approx.) is 32 mA.}
\end{qoptions}
\sect{SECTION-I(ii)}
\qn{9} A pair of parallel rails of negligible resistance is fixed at a distance $\ell=10$ cm on the slope of an inclined plane making an angle $\alpha=37^\circ$ with the horizontal. A battery of EMF $U_0$ is connected across the lower ends of the rails. A horizontal rod of mass $M=1$kg and resistance 1 m$\Omega$ is placed perpendicularly to the rails. The rod can move without any friction. The rod is connected to a block of mass $m=2$kg through a light thread as shown in diagram. The arrangement is placed in a vertically downward uniform magnetic field $B=1$T. If $U_0$ is equal to zero :
\qfig{p2_phy_s2_q09_1}
{\footnotesize\setlength{\tabcolsep}{2pt}\begin{tabular}{|c|p{0.27\linewidth}|c|p{0.14\linewidth}|}
\hline
 & \centering\textbf{List--I} & & \textbf{List--II}\\ \hline
(I) & Initial acceleration of the blocks (in m/s$^2$) is & (P) & 0\\ \hline
(II) & Initial current through the resistance (in A) is & (Q) & $\dfrac{14}{3}$\\ \hline
(III) & Terminal speed attained by the rod (in m/s) is & (R) & $\dfrac{245}{8}$\\ \hline
(IV) & Power loss in resistance (in W) when rod moves at terminal velocity. & (S) & 175\\ \hline
 & & (T) & $\dfrac{35}{16}$\\ \hline
 & & (U) & $\dfrac{7840}{9}$\\ \hline
\end{tabular}}\par
\begin{qoptions}
\opt{A}{I $\to$ Q;II $\to$ P;III $\to$ T;IV $\to$ R}
\opt{B}{I $\to$ P;II $\to$ Q;III $\to$ T;IV $\to$ S}
\opt{C}{I $\to$ R;II $\to$ T;III $\to$ P;IV $\to$ S}
\opt{D}{I $\to$ Q;II $\to$ P;III $\to$ R;IV $\to$ T}
\end{qoptions}
\qn{10} A pair of parallel rails of negligible resistance is fixed at a distance $\ell=10$ cm on the slope of an inclined plane making an angle $\alpha=37^\circ$ with the horizontal. A battery of EMF $U_0$ is connected across the lower ends of the rails. A horizontal rod of mass $M=1$kg and resistance 1 m$\Omega$ is placed perpendicularly to the rails. The rod can move without any friction. The rod is connected to a block of mass $m=2$kg through a light thread as shown in diagram. The arrangement is placed in a vertically downward uniform magnetic field $B=1$T. If $U_0$ is equal to 0.175V :
\qfig{p2_phy_s2_q10_1}
{\footnotesize\setlength{\tabcolsep}{2pt}\begin{tabular}{|c|p{0.27\linewidth}|c|p{0.14\linewidth}|}
\hline
 & \centering\textbf{List--I} & & \textbf{List--II}\\ \hline
(I) & Initial acceleration of the blocks (in m/s$^2$) is & (P) & 0\\ \hline
(II) & Initial current through the resistance (in A) is & (Q) & $\dfrac{14}{3}$\\ \hline
(III) & Terminal speed attained by the rod (in m/s) is & (R) & $\dfrac{245}{8}$\\ \hline
(IV) & Power loss in resistance (in W) when rod attains terminal speed is & (S) & 175\\ \hline
 & & (T) & $\dfrac{35}{16}$\\ \hline
 & & (U) & $\dfrac{7840}{9}$\\ \hline
\end{tabular}}\par
\begin{qoptions}
\opt{A}{I $\to$ Q;II $\to$ S;III $\to$ P;IV $\to$ R}
\opt{B}{I $\to$ P;II $\to$ S;III $\to$ P;IV $\to$ R}
\opt{C}{I $\to$ P;II $\to$ Q;III $\to$ R;IV $\to$ S}
\opt{D}{I $\to$ P;II $\to$ R;III $\to$ P;IV $\to$ Q}
\end{qoptions}
\qn{11} A cylindrical isotropic solid of coefficient of linear expansion $\alpha$ and density $\rho$ floats in a liquid of coefficient of volume expansion $\gamma$ and density $d$ as shown in the figure.
\qfig{p2_phy_s2_q11_1}
{\footnotesize\setlength{\tabcolsep}{2pt}\begin{tabular}{|c|p{0.27\linewidth}|c|p{0.14\linewidth}|}
\hline
 & \centering\textbf{List--I} & & \textbf{List--II}\\ \hline
(I) & Volume of cylinder inside the liquid remain constant & (P) & $\gamma=0$\\ \hline
(II) & Volume of cylinder outside the liquid remain constant. & (Q) & $\gamma=2\alpha$\\ \hline
(III) & Height of cylinder outside the liquid remains constant & (R) & $\gamma=3\alpha$\\ \hline
(IV) & Height of cylinder inside the liquid remain constant & (S) & $\gamma=4\alpha$\\ \hline
 & & (T) & $\gamma=5\alpha$\\ \hline
 & & (U) & $\gamma=6\alpha$\\ \hline
\end{tabular}}\par
If $\dfrac\rho d=1$ then
\begin{qoptions}
\opt{A}{I $\to$ P;II $\to$ R;III $\to$ R;IV $\to$ U}
\opt{B}{I $\to$ P;II $\to$ R;III $\to$ R;IV $\to$ Q}
\opt{C}{I $\to$ P;II $\to$ S;III $\to$ U;IV $\to$ T}
\opt{D}{I $\to$ Q;II $\to$ S;III $\to$ T;IV $\to$ U}
\end{qoptions}
\qn{12} A cylindrical isotropic solid of coefficient of linear expansion $\alpha$ and density $\rho$ floats in a liquid of coefficient of volume expansion $\gamma$ and density $d$ as shown in the figure.
\qfig{p2_phy_s2_q12_1}
{\footnotesize\setlength{\tabcolsep}{2pt}\begin{tabular}{|c|p{0.27\linewidth}|c|p{0.14\linewidth}|}
\hline
 & \centering\textbf{List--I} & & \textbf{List--II}\\ \hline
(I) & Volume of cylinder inside the liquid remain constant & (P) & $\gamma=0$\\ \hline
(II) & Volume of cylinder outside the liquid remain constant. & (Q) & $\gamma=2\alpha$\\ \hline
(III) & Height of cylinder outside the liquid remains constant & (R) & $\gamma=3\alpha$\\ \hline
(IV) & Height of cylinder inside the liquid remain constant & (S) & $\gamma=4\alpha$\\ \hline
 & & (T) & $\gamma=5\alpha$\\ \hline
 & & (U) & $\gamma=6\alpha$\\ \hline
\end{tabular}}\par
If $\dfrac\rho d=\dfrac12$ then
\begin{qoptions}
\opt{A}{I $\to$ Q;II $\to$ U;III $\to$ S;IV $\to$ P}
\opt{B}{I $\to$ P;II $\to$ U;III $\to$ S;IV $\to$ Q}
\opt{C}{I $\to$ P;II $\to$ U;III $\to$ T;IV $\to$ S}
\opt{D}{I $\to$ P;II $\to$ T;III $\to$ U;IV $\to$ Q}
\end{qoptions}
\sect{SECTION-II}
\qn{1} A conductor $OP$ of length $\ell$ placed in $x$-$y$ plane with one end hinged at origin. In this region a non-uniform magnetic field exist along positive $z$ direction of which magnitude depends upon its $y$-coordinate which is given as $B=B_0\left(1+\dfrac{y^2}{\ell^2}\right)$ tesla. If conductor $OP$ start rotating with angular velocity $-\omega_0\hat k$. Then the potential difference across $OP$ is $N.B_0\omega\ell^2$ when the conductor lies along $y$-axis. Find $N$.
\qfig{p2_phy_s3_q01_1}
\qn{2} In series RLC circuit with $20\Omega$ resistance is connected an AC source. When only the capacitance is removed the current leads the voltage by $53^\circ$ and if only inductor is removed voltage leads by $37^\circ$. The impedence of the circuit is $N\ \Omega$. Find the value of $N$.[Take $\tan37^\circ=3/4$ and $\sqrt{193}=13.89$]
\qn{3} The U-shaped tube shown in figure is filled with water. Now current $I$ is made to flow in the solenoid of length $\ell$ and number of turns $N$, tightly wrapped around the left arm of the tube. Assume that diameter of tube is very less as compared to the length of the solenoid and the relative permeability of water $\mu_r$ is the change in the level of liquid is $\dfrac{\alpha\mu_0N^2I^2(1-\mu_r)}{\ell^2\rho g}$. The value of $\alpha$ is.
\qfig{p2_phy_s3_q03_1}
\qn{4} A heating element is used to raise the temperature of water of mass 10 kg inside the boiler. The heater is kept on for a very long time and it is observed that no steam is formed. Now the heating element is switched off and temperature of water is observed to decrease at a rate of $1^\circ$C/min. The water equivalent of boiler is 3 kg. The power (in W) rating of heating element is (Take specific heat capacity of water 4200 J/kg-K)
\qn{5} A uniform cube is kept on a horizontal surface. What is the ratio of strain energy stored in lower one third to the upper one third of cube.
\qn{6} In the figure shown, the hollow cone of mass $m$ and radius $R$ is free to rotate about axis $OO'$. Its curved surface carries a uniform surface charge density $\sigma$ as shown. The material of the cone is non conducting. If a time varying magnetic field, $\vec B=\vec kt$ is switched on at $t=0$ in the region, where the vector $\vec k$ is along $OO'$, the angular velocity `$\omega$' acquired by the cone due to the electromagnetic forces $\dfrac{Nk\sigma\pi t}{6m}$, Find $N$.
\qfig{p2_phy_s3_q06_1}
\partsec{2}{CHEMISTRY}
\sect{SECTION-I(i)}
\qn{1} In order to study effect of temperature on rate of reaction, another term ``$\sigma$'' is defined as $\sigma=\dfrac1K\dfrac{dK}{dT}$ [where $K$ is rate constant and $T$ represents temperature]. If for a reaction value of $\sigma=\dfrac{1.25\times10^6}{T^3}$ [$T$ should be in Kelvin], then identify the correct statement(s).
\begin{qoptions}
\opt{A}{Activation energy for the given reaction at 250 Kelvin is 10 Kcal/mol.}
\opt{B}{Activation energy for the given reaction at 2000 Kelvin is 1.25 Kcal/mol.}
\opt{C}{Rate of increase of rate constant with temperature is higher at lower temperature range than at higher temperature range.}
\opt{D}{The value of $\dfrac{dK}{dT}$ at 2000 Kelvin is $0.625\times K$.}
\end{qoptions}
\qn{2} Select the correct statement(s) :
\begin{qoptions}
\opt{A}{$\mathrm{C_{Vm}}$ is state function}
\opt{B}{Depression in freezing point ($\Delta T_f$) is state variable}
\opt{C}{A process can be defined completely by knowing initial and final states of system.}
\opt{D}{Poission's ratio ($\gamma$) for $\mathrm{NH_3(g)}$ will be $\dfrac{11}{9}$ if 25\% vibration degree of freedoms are activated.}
\end{qoptions}
\qn{3} The complex $\mathrm{K_4[Zn(CN)_4(O_2)_2]}$ is oxidised to $\mathrm{K_2[Zn(CN)_4(O_2)_2]}$, then. Which of the following is/are correct ?
\begin{qoptions}
\opt{A}{Zn(II) is oxidised to Zn(IV)}
\opt{B}{O-O bond length decreases}
\opt{C}{Magnetic moment increases}
\opt{D}{Magnetic behaviour does not change}
\end{qoptions}
\qn{4} The correct option(s) regarding the complex $\mathrm{[Co(en)(NH_3)_3(H_2O)]^{3+}}$ :\\
(en = $\mathrm{H_2NCH_2CH_2NH_2}$) is (are)
\begin{qoptions}
\opt{A}{It has two geometrical isomers}
\opt{B}{It will have three geometrical isomers if bidentate `en' is replaced by two cyanide ligands}
\opt{C}{It is paramagnetic}
\opt{D}{It absorbs light at longer wavelength as compared to $\mathrm{[Co(en)(NH_3)_4]^{3+}}$}
\end{qoptions}
\qn{5} Choose the \textbf{CORRECT} order for the given properties
\begin{qoptions}
\opt{A}{$\mathrm{MgSO_4<SrSO_4<BaSO_4}$ : Thermal stability order}
\opt{B}{$\mathrm{BeC_2O_4<CaC_2O_4<BaC_2O_4}$ : Solubility in water}
\opt{C}{$\mathrm{LiCl>NaCl>KCl}$ : Melting point order}
\opt{D}{$\mathrm{BeF_2>CaF_2>SrF_2}$ : Covalent character order}
\end{qoptions}
\qn{6} Consider the following reaction\par
$\mathrm{\left[Co\left(\eta^5-C_5H_5\right)(CO)_2\right]}$ \textsubscript{(Compound I)} $+\ \mathrm{C_3H_5I}\longrightarrow\mathrm{\left[Co\left(\eta^3-C_3H_5\right)\left(\eta^5-C_5H_5\right)I\right]}$ \textsubscript{(Compound II)} $+\ 2\mathrm{CO}$\par
Which of the following statement is incorrect regarding the above reaction?
\begin{qoptions}
\opt{A}{I and II both consists of trihapto ligands}
\opt{B}{The above reaction is a redox reaction}
\opt{C}{Here CO acts as 3 electron donor for complex-I}
\opt{D}{Both complexes are associated with synergic bonding}
\end{qoptions}
\qn{7} \qfig{p2_chem_s1_q07_1}
Which of the following are correct?
\begin{qoptions}
\opt{A}{(P) is \qfigi{p2_chem_s1_q07_2}}
\opt{B}{(Q) is \qfigi{p2_chem_s1_q07_3}}
\opt{C}{(P) is \qfigi{p2_chem_s1_q07_4}}
\opt{D}{(Q) is \qfigi{p2_chem_s1_q07_5}}
\end{qoptions}
\qn{8} Which of the following reactions are correct?
\begin{qoptions}
\opt{A}{\qfigi{p2_chem_s1_q08_A} $\xrightarrow{\mathrm{LiAlH_4}}$ \qfigi{p2_chem_s1_q08_2}}
\opt{B}{\qfigi{p2_chem_s1_q08_B} \qfigi{p2_chem_s1_q08_4}}
\opt{C}{\qfigi{p2_chem_s1_q08_C} $\xrightarrow{\mathrm{NaBH_4}}$ \qfigi{p2_chem_s1_q08_6}}
\opt{D}{$\mathrm{CH_3{-}C{\equiv}N\xrightarrow{LiAlH_4}CH_3{-}CH_2{-}NH_2}$}
\end{qoptions}
\sect{SECTION-I(ii)}
\qn{9} Match \textbf{List-I} with \textbf{List-II} and choose the correct option :\par
{\footnotesize\setlength{\tabcolsep}{2pt}\begin{tabular}{|c|c|c|p{0.38\linewidth}|}
\hline
\multicolumn{2}{|c|}{\textbf{List-I}} & \multicolumn{2}{c|}{\textbf{List-II}}\\ \hline
(I) & $Z=\dfrac38$ & (P) & At Boyle's temperature in low to moderate pressure region.\\ \hline
(II) & $Z>1$ & (Q) & At inversion temperature\\ \hline
(III) & $Z<1$ & (R) & $\mathrm{H_2}$(gas), $P\to0$, $T=0^\circ$C\\ \hline
(IV) & $Z=1$ & (S) & $\mathrm{H_2}$(gas), $P=1$ bar, $T=0^\circ$C\\ \hline
 & & (T) & $\mathrm{CO_2}$(gas), $P=1$ bar, $T=0^\circ$C\\ \hline
 & & (U) & $\mathrm{CO_2}$ at critical point.\\ \hline
\end{tabular}}
\begin{qoptions}
\opt{A}{I $\to$ Q;II $\to$ Q;III $\to$ T;IV $\to$ P}
\opt{B}{I $\to$ U;II $\to$ S;III $\to$ U;IV $\to$ T}
\opt{C}{I $\to$ U;II $\to$ Q;III $\to$ T;IV $\to$ P,R}
\opt{D}{I $\to$ Q;II $\to$ Q;III $\to$ T;IV $\to$ U}
\end{qoptions}
\qn{10} Match \textbf{List-I} with \textbf{List-II} and choose the incorrect option :\par
{\footnotesize\setlength{\tabcolsep}{2pt}\begin{tabular}{|c|c|c|p{0.38\linewidth}|}
\hline
\multicolumn{2}{|c|}{\textbf{List-I}} & \multicolumn{2}{c|}{\textbf{List-II}}\\ \hline
(I) & $Z=\dfrac38$ & (P) & At Boyle's temperature in low to moderate pressure region.\\ \hline
(II) & $Z>1$ & (Q) & At inversion temperature\\ \hline
(III) & $Z<1$ & (R) & $\mathrm{H_2}$(gas), $P\to0$, $T=0^\circ$C\\ \hline
(IV) & $Z=1$ & (S) & $\mathrm{H_2}$(gas), $P=1$ bar, $T=0^\circ$C\\ \hline
 & & (T) & $\mathrm{CO_2}$(gas), $P=1$ bar, $T=0^\circ$C\\ \hline
 & & (U) & $\mathrm{CO_2}$ at critical point.\\ \hline
\end{tabular}}
\begin{qoptions}
\opt{A}{I $\to$ U;II $\to$ S;III $\to$ U;IV $\to$ P}
\opt{B}{I $\to$ U;II $\to$ Q;III $\to$ U;IV $\to$ P}
\opt{C}{I $\to$ U;II $\to$ Q,S;III $\to$ T,U;IV $\to$ P,R}
\opt{D}{I $\to$ U;II $\to$ S;III $\to$ T;IV $\to$ S}
\end{qoptions}
\qn{11} {\footnotesize\setlength{\tabcolsep}{2pt}\begin{tabular}{|c|p{0.3\linewidth}|c|p{0.22\linewidth}|}
\hline
 & \centering\textbf{List-I} & & \textbf{List-II}\\ \hline
(I) & $(-)$-1-Bromo-2-ethyl pentane (Single enantiomer) $\xrightarrow[\mathrm{S_N2}]{\mathrm{aq.NaOH}}$ & (P) & Inversion of configuration\\ \hline
(II) & $(-)$-2-Bromopentane (single enantiomer) $\xrightarrow[\mathrm{S_N2}]{\mathrm{aq.NaOH}}$ & (Q) & Retention of configuration\\ \hline
(III) & $(-)$-3-Bromo-3-methyl hexane (single enantiomer) $\xrightarrow[\mathrm{S_N2}]{\mathrm{aq.NaOH}}$ & (R) & Mixture of enantiomers\\ \hline
(IV) & \qfigi{p2_chem_s2_q11_1} & (S) & Mixture of structural isomers\\ \hline
 & & (T) & Mixture of diastereomers\\ \hline
 & & (U) & Mixture of ring chain isomers\\ \hline
\end{tabular}}\par
Which of the following is correct?
\begin{qoptions}
\opt{A}{I $\to$ P}
\opt{B}{II $\to$ Q}
\opt{C}{III $\to$ S}
\opt{D}{IV $\to$ T}
\end{qoptions}
\qn{12} {\footnotesize\setlength{\tabcolsep}{2pt}\begin{tabular}{|c|p{0.3\linewidth}|c|p{0.22\linewidth}|}
\hline
 & \centering\textbf{List-I} & & \textbf{List-II}\\ \hline
(I) & $(-)$-1-Bromo-2-ethyl pentane (Single enantiomer) $\xrightarrow[\mathrm{S_N2}]{\mathrm{aq.NaOH}}$ & (P) & Inversion of configuration\\ \hline
(II) & $(-)$-2-Bromopentane (single enantiomer) $\xrightarrow[\mathrm{S_N2}]{\mathrm{aq.NaOH}}$ & (Q) & Retention of configuration\\ \hline
(III) & $(-)$-3-Bromo-3-methyl hexane (single enantiomer) $\xrightarrow[\mathrm{S_N2}]{\mathrm{aq.NaOH}}$ & (R) & Mixture of enantiomers\\ \hline
(IV) & \qfigi{p2_chem_s2_q11_2} & (S) & Mixture of structural isomers\\ \hline
 & & (T) & Mixture of diastereomers\\ \hline
 & & (U) & Mixture of ring chain isomers\\ \hline
\end{tabular}}\par
Which of the following is incorrect?
\begin{qoptions}
\opt{A}{I $\to$ Q}
\opt{B}{II $\to$ P}
\opt{C}{III $\to$ S}
\opt{D}{IV $\to$ T}
\end{qoptions}
\sect{SECTION-II}
\qn{1} A 15\% by mass of sucrose solution has undergone partial inversion\par
Sucrose (aq) + Water $\longrightarrow$ Glucose(aq) + Fructose (aq)\par
After 10 hours from the start of reaction the elevation in boiling point of solution is $\dfrac{103}{200}$ K\par
The reaction follows $1^{\mathrm{st}}$ order Kinetics with respect to Sucrose. Find time (in hours) for 87.5\% of inversion of sucrose ? Neglect the change in mass of water due to reaction.\par
(Given : $K_b$(water) = 0.57 kg K mol$^{-1}$, Molar mass of sucrose = 342 gm/mol)
\qn{2} For hypothetical reaction occurring in same container\par
$\mathrm{A+8B\longrightarrow P}$\quad $\Delta H=-2x$ KJ\par
$\mathrm{M\longrightarrow 2Q+R}$\quad $\Delta H=x$ KJ\par
Temperature of solution is increased by $2^\circ$C per sec.\par
If rate of disappearance of B = 2 mole sec$^{-1}$\par
Heat capacity of overall reaction mixture = 100 x J/K then what is rate of appearance of Q (in mol/sec).
\qn{3} Find the sum of number of unpaired electrons for following complexes\par
$\mathrm{[PtCl_4]^{-2}}$, $\mathrm{[Pt(NH_3)(py)ClBr]}$, $\mathrm{[Cr(C_2O_4)_3]^{-3}}$, $\mathrm{[Co(OX)_3]^{-3}}$, $\mathrm{[FeF_6]^{-3}}$, $\mathrm{[NiCl_4]^{-2}}$
\qn{4} Among the following, total number of specie(s) is/are exist as optically inactive form :\par
$\mathrm{[Fe(en)_3]Cl_3}$ ; $\mathrm{[RhCl(PPh_3)(CO)(H_2O)]}$ ; $\mathrm{[Pd(en)Cl_2]}$ ; \qfigi{p2_chem_s3_q04_1}
\qfullbegin
\qn{5} How many organic substrates will show rearrangement during $\mathrm{S_N1}$ reaction?
\qfig{p2_chem_s3_q05_1}
\qfullsplit
\qfig{p2_chem_s3_q05_2}
\qfig{p2_chem_s3_q05_3}
\qfullend
\qn{6} $\mathrm{Me-\overset{OH}{\overset{|}{CH}}-CH{=}CH-Me}$\par
Total number of stereoisomers?
\partsec{3}{MATHEMATICS}
\sect{SECTION-I(i)}
\qn{1} Let $f(x,y)=\dfrac{1}{\displaystyle\lim_{\alpha\to0}\left(\dfrac{\sqrt[x]{\sec\alpha}-\sqrt[y]{\sec\alpha}}{\alpha^2}\right)}$, then
\begin{qoptions}
\opt{A}{$f(2,3)=12$}
\opt{B}{$f(x,y)=12$ has infinitely many solutions}
\opt{C}{Slope of tangent to the curve $f(x,y)=12$ at $x=2$ is 4}
\opt{D}{$f(x,y)=12$ has no horizontal tangent}
\end{qoptions}
\qn{2} If $\displaystyle\int\cot\left(2\tan^{-1}\sqrt{\dfrac{\sqrt{1+\sqrt x}-x^{1/4}}{\sqrt{1+\sqrt x}+x^{1/4}}}\right)dx=\dfrac{qx^{p/q}}{p}+C$; (where $p$ and $q$ are relatively prime and $C$ is constant of integration), then
\begin{qoptions}
\opt{A}{$p=5$}
\opt{B}{$p+2q=13$}
\opt{C}{$p-q=1$}
\opt{D}{$q=5$}
\end{qoptions}
\qn{3} If $y^x-x^y=1$, then $\dfrac{dy}{dx}$ at $x=1$, is -
\begin{qoptions}
\opt{A}{$2(1-\ell n2)$}
\opt{B}{$2(1+\ell n2)$}
\opt{C}{$\ell n(4e^2)$}
\opt{D}{$\ell n\left(\dfrac{e^2}{4}\right)$}
\end{qoptions}
\qn{4} For the equation $\begin{vmatrix}\tan x&\cot x&\cot x\\\cot x&\tan x&\tan x\\\cot x&\cot x&\tan x\end{vmatrix}=0$ which of the following is/are incorrect ?
\begin{qoptions}
\opt{A}{These are four distinct real roots in $[0,2\pi]$}
\opt{B}{$\tan x=\sqrt2$}
\opt{C}{These are six distinct real roots in $[0,2\pi]$}
\opt{D}{Number of distinct real roots in $\left[-\dfrac\pi2,\dfrac\pi2\right]$ is four}
\end{qoptions}
\qn{5} If $\displaystyle\int\dfrac{x+\sqrt[3]{x^2}+\sqrt[6]{x}}{x\left(1+\sqrt[3]{x}\right)}dx=A.x^B+C\tan^{-1}\left(x^D\right)+E$, then which of the following is/are true (where `E' is constant of integration) ?
\begin{qoptions}
\opt{A}{$A=\dfrac23$}
\opt{B}{$AB=1$}
\opt{C}{$(AB+C)$ is a prime number}
\opt{D}{Sum of all coefficients of even power of $x$ in the expansion of $\left(AB+\dfrac xD\right)^C$ is equal to $\dfrac{(C+1)^C+(C-1)^C}{2}$}
\end{qoptions}
\qn{6} Let $f:R\to A$ be an invertible function defined by $f(x)=\log_2\left(2^{x^3+x}+4\right)$. If $g(x)$ be its inverse function, then -
\begin{qoptions}
\opt{A}{$g(3)=2$}
\opt{B}{$g'(3)=\dfrac12$}
\opt{C}{$\displaystyle\lim_{x\to3}(g(x))^{\frac{\ell n4}{x-3}}=2$}
\opt{D}{set $A$ is $(2,\infty)$}
\end{qoptions}
\qn{7} If the function $f(x)=-4e^{\frac{1-x}{2}}+1+x+\dfrac{x^2}{2}+\dfrac{x^3}{3}$ and $g(x)=f^{-1}(x)$, then -
\begin{qoptions}
\opt{A}{$g'\left(-\dfrac76\right)=\dfrac15$}
\opt{B}{$g'\left(\dfrac76\right)=\dfrac15$}
\opt{C}{$g(1)=-\dfrac76$}
\opt{D}{$g\left(-\dfrac76\right)=1$}
\end{qoptions}
\qn{8} For the figure given below if $\tan\alpha=\dfrac mn$, where $m$ and $n$ are relatively prime positive integers then
\qfig{p2_math_s1_q08_1}
\begin{qoptions}
\opt{A}{$m+n=463$}
\opt{B}{$n-m=122$}
\opt{C}{$2m-n=41$}
\opt{D}{$m+n=465$}
\end{qoptions}
\sect{SECTION-I(ii)}
\qn{9} If a curve $f(x,y,a,b)=0$ is given by $\sin^{-1}x+\cos^{-1}ay+\cos^{-1}(bxy)=\dfrac\pi2$\par
Match List-I with List-II and select the correct answer using the code given below the list.\par
{\footnotesize\setlength{\tabcolsep}{2pt}\begin{tabular}{|c|p{0.3\linewidth}|c|p{0.3\linewidth}|}
\hline
\multicolumn{2}{|c|}{\textbf{List-I}} & \multicolumn{2}{c|}{\textbf{List-II}}\\ \hline
(I) & if $a=1$, $b=2$, then & (P) & $(1-x^2)(1-y^2)=x^2y^2$\\ \hline
(II) & if $a=2$, $b=1$, then & (Q) & $a+b=8$\\ \hline
(III) & if $a=0$, $b=1$, then & (R) & $\dfrac{dy}{dx}=-\left(\dfrac{y^2+1}{xy}\right)$, $x,y\in R_0$\\ \hline
(IV) & if $a=b$ and area bounded by rectangle formed by $f(x,y,a,b)=0$ is 1 sq. unit, then & (S) & $x^2+4y^2=1+3x^2y^2$\\ \hline
 & & (T) & $x+y=2$\\ \hline
 & & (U) & $x^2+y^2=3x^2y^2$\\ \hline
\end{tabular}}\par
Which of the following is the only \textbf{CORRECT} combination ?
\begin{qoptions}
\opt{A}{I $\to$ R}
\opt{B}{II $\to$ Q}
\opt{C}{III $\to$ R}
\opt{D}{IV $\to$ S}
\end{qoptions}
\qn{10} If a curve $f(x,y,a,b)=0$ is given by $\sin^{-1}x+\cos^{-1}ay+\cos^{-1}(bxy)=\dfrac\pi2$\par
Match List-I with List-II and select the correct answer using the code given below the list.\par
{\footnotesize\setlength{\tabcolsep}{2pt}\begin{tabular}{|c|p{0.3\linewidth}|c|p{0.3\linewidth}|}
\hline
\multicolumn{2}{|c|}{\textbf{List-I}} & \multicolumn{2}{c|}{\textbf{List-II}}\\ \hline
(I) & if $a=1$, $b=2$, then & (P) & $(1-x^2)(1-y^2)=x^2y^2$\\ \hline
(II) & if $a=2$, $b=1$, then & (Q) & $a+b=8$\\ \hline
(III) & if $a=0$, $b=1$, then & (R) & $\dfrac{dy}{dx}=-\left(\dfrac{y^2+1}{xy}\right)$, $x,y\in R_0$\\ \hline
(IV) & if $a=b$ and area bounded by rectangle formed by $f(x,y,a,b)=0$ is 1 sq. unit, then & (S) & $x^2+4y^2=1+3x^2y^2$\\ \hline
 & & (T) & $x+y=2$\\ \hline
 & & (U) & $x^2+y^2=3x^2y^2$\\ \hline
\end{tabular}}\par
Which of the following is the only \textbf{INCORRECT} combination ?
\begin{qoptions}
\opt{A}{I $\to$ P}
\opt{B}{II $\to$ S}
\opt{C}{III $\to$ R}
\opt{D}{IV $\to$ R}
\end{qoptions}
\qn{11} Match List-I with List-II and select the correct answer using the code given below the list.\par
{\footnotesize\setlength{\tabcolsep}{2pt}\begin{tabular}{|c|p{0.27\linewidth}|c|p{0.1\linewidth}|}
\hline
 & \centering\textbf{List-I} & & \textbf{List-II}\\ \hline
(I) & Number of solutions of $4(\sin^6\theta+\cos^6\theta)=10-18\sin\theta\cos\theta$ in $[0,3\pi]$, is & (P) & 2\\ \hline
(II) & Number of possible pairs $(x,y)$ satisfying $x^2-6x+10-\sin^2y=0$, where $y\in[0,2\pi]$ and $x\in R$, is & (Q) & 3\\ \hline
(III) & Number of integral values of $x$ satisfying $\sin x<\dfrac12$ in $[0,2\pi]$, is & (R) & 4\\ \hline
(IV) & Number of solutions satisfying $\tan x+\sec x=2\cos x$ in $[0,4\pi]$, is & (S) & 5\\ \hline
 & & (T) & 6\\ \hline
 & & (U) & 7\\ \hline
\end{tabular}}\par
Which of the following is the only \textbf{CORRECT} combination ?
\begin{qoptions}
\opt{A}{I $\to$ R}
\opt{B}{II $\to$ P}
\opt{C}{III $\to$ R}
\opt{D}{IV $\to$ Q}
\end{qoptions}
\qn{12} Match List-I with List-II and select the correct answer using the code given below the list.\par
{\footnotesize\setlength{\tabcolsep}{2pt}\begin{tabular}{|c|p{0.27\linewidth}|c|p{0.1\linewidth}|}
\hline
 & \centering\textbf{List-I} & & \textbf{List-II}\\ \hline
(I) & Number of solutions of $4(\sin^6\theta+\cos^6\theta)=10-18\sin\theta\cos\theta$ in $[0,3\pi]$, is & (P) & 2\\ \hline
(II) & Number of possible pairs $(x,y)$ satisfying $x^2-6x+10-\sin^2y=0$, where $y\in[0,2\pi]$ and $x\in R$, is & (Q) & 3\\ \hline
(III) & Number of integral values of $x$ satisfying $\sin x<\dfrac12$ in $[0,2\pi]$, is & (R) & 4\\ \hline
(IV) & Number of solutions satisfying $\tan x+\sec x=2\cos x$ in $[0,4\pi]$, is & (S) & 5\\ \hline
 & & (T) & 6\\ \hline
 & & (U) & 7\\ \hline
\end{tabular}}\par
Which of the following is the only \textbf{INCORRECT} combination ?
\begin{qoptions}
\opt{A}{I $\to$ Q}
\opt{B}{II $\to$ P}
\opt{C}{III $\to$ S}
\opt{D}{IV $\to$ S}
\end{qoptions}
\sect{SECTION-II}
\qn{1} Let $J=\displaystyle\int\dfrac{x+y}{xy}dy$ and $I=\displaystyle\int\dfrac{x+y}{x^2}dx$ where $x$ and $y$ are independent variables such that $g\left(\dfrac yx\right)=J-I$ and $g(1)=1$, then value of $g(e)$ is
\qn{2} Let $f(x)$ be a differentiable function on $[0,8]$ such that $f(1)=5$, $f(2)=\dfrac12$, $f(3)=6$, $f(4)=-3$, $f(5)=4$, $f(6)=\dfrac14$ and $f(7)=\dfrac{-1}{5}$. If the minimum number of roots of the equation $f'(x)-f'(x)(f(x))^2=0$ is $\lambda$, then $\dfrac\lambda3$ equals
\qn{3} Let $f(x)$ be a real valued function satisfying $2f(xy)=(f(x))^y+(f(y))^x\ \forall x,y\in R$, $f(1)=3$ and $\displaystyle\lim_{x\to0}\dfrac{f(x)+f(-x)-f(f(0))+1}{x^2}$ is equal to $(\ell nk)^2$ where $k$ is natural number, then the value of $\dfrac k2$ is
\qn{4} $\displaystyle\lim_{x\to\infty}x^3\left(\sqrt{x^2+\sqrt{1+x^4}}-\sqrt2x\right)=\dfrac{\sqrt2}{\lambda}$ find value of $\lambda$
\qn{5} Let $f:(0,\infty)\to(0,\infty)$ be a function differentiable at 3 and satisfying $f(3)=3f'(3)$. If the value of $\displaystyle\lim_{x\to\infty}\left(\dfrac{f\left(3+\frac3x\right)}{f(3)}\right)^x$ is $L$ then $\dfrac{[L]}{5}$ is (where $[.]$ is greatest integer function)
\qn{6} If a rectangular tank has a sliding panel $S$ that divides it into two adjustable tanks of width 3 ft. See figure. Water is poured into the left compartment at the rate of 5 ft$^3$/min. At the same time $S$ is moved to the right at the rate of 3ft/min. When the left compartment is 10 ft long it contains 70ft$^3$ of water, then how fast the water level is falling
\qfig{p2_math_s3_q06_1}
\end{multicols}
\end{document}
